% --- Archon physics formalization source begin ---
% archon:physics
% archon:covers PhyXMiniProblems/problem_phyx_mini_0891.lean
% archon:source-report reports/phyx_mini/problem_phyx_mini_0891.source.json

\chapter{Physics problem phyx\_mini\_0891}
\label{ch:PhyXMiniProblems_problem_phyx_mini_0891}

\paragraph{Problem source.}
\#\# Physical scenario

A horizontal electric field causes the charged ball in Figure to hang at a 15\textasciicircum{}\textbackslash{}circ angle, as shown. The spring is plastic, so it doesn’t discharge the ball, and in its equilibrium position the spring extends only to the vertical dashed line.

\#\# Question

What is the electric field strength?

\#\# Answer choices

- A: A: 1.25 \textbackslash{}times 10\textasciicircum{}\{3\}N/C

- B: B: 1.35 \textbackslash{}times 10\textasciicircum{}\{5\}N/C

- C: C: 7.5 \textbackslash{}times 10\textasciicircum{}\{5\}N/C

- D: D: 5.33 \textbackslash{}times 10\textasciicircum{}\{5\}N/C

\#\# Figure caption (auxiliary; use the image as primary evidence)

The image depicts a physics diagram involving a spring, a charged object, and an electric field. Here's a detailed description:

- **Spring**: Attached to a wall on the left side, with a spring constant of 0.050 N/m.
- **Object**: A red sphere labeled with a mass of 3.0 grams and a charge of 20 nC (nanocoulombs).
- **Displacement**: The spring is stretched to the right, and the object is positioned at the end of the spring, 60 cm from the wall.
- **Angle**: The line from the wall to the object is shown at a 15-degree angle from the vertical.
- **Electric Field**: Denoted by \textbackslash{}( \textbackslash{}mathbf\{E\} \textbackslash{}), with arrows indicating the direction to the right across the image.
- **Text and Labels**: 
  - Spring constant is labeled as 0.050 N/m.
  - Object's weight and charge are labeled as 3.0 g and 20 nC, respectively.
  - The length from the wall to the object along the spring is marked as 60 cm.

This setup likely relates to problems involving mechanics and electromagnetism, such as balancing forces on a charged object in an electric field.

\#\# Dataset metadata

Electrostatics; Physical Model Grounding Reasoning, Multi-Formula Reasoning

\paragraph{Recorded answer/context.}
C: C: 7.5 \textbackslash{}times 10\textasciicircum{}\{5\}N/C

\paragraph{Figure/image path.}
/root/proposal\_for\_physic/science-mango/phyx\_mini\_run/phyx\_data/test\_image/891.png

\paragraph{Formalization target.}
create a compiling Lean file with sorry bodies at `PhyXMiniProblems/problem\_phyx\_mini\_0891.lean`. The Lean declarations must preserve the physical quantities, dimensions or dimensional roles, figure labels, governing-law hypotheses, and final relation expressed by this problem.
Use Mathlib/Physlib names found through LeanExplore where available. If a physics API is missing, introduce faithful local abstractions rather than scalar placeholder aliases.

\begin{theorem}[Physics formalization target]
\label{thm:physics:phyx_mini_0891:target}
\lean{PhyXMiniProblems.ProblemPhyXMini0891.electricFieldStrength_agreesWithChoiceC}
\uses{def:physics:phyx-mini-0891:phyxminiproblems-problemphyxmini0891-chargedballspringsetup, def:physics:phyx-mini-0891:phyxminiproblems-problemphyxmini0891-matcheschargedballspringfigure, def:physics:phyx-mini-0891:phyxminiproblems-problemphyxmini0891-matcheschargedballspringscenario, def:physics:phyx-mini-0891:phyxminiproblems-problemphyxmini0891-usesstandardterrestrialgravity, def:physics:phyx-mini-0891:phyxminiproblems-problemphyxmini0891-hasphysicalchargedballspringparameters, def:physics:phyx-mini-0891:phyxminiproblems-problemphyxmini0891-satisfiesdeflectedtetherandspringgeometry, def:physics:phyx-mini-0891:phyxminiproblems-problemphyxmini0891-satisfiesuniformhorizontalelectricfield, def:physics:phyx-mini-0891:phyxminiproblems-problemphyxmini0891-satisfieschargedballstaticequilibrium, def:physics:phyx-mini-0891:phyxminiproblems-problemphyxmini0891-isuniqueclosestanswer}
The assigned autoformalize agent should translate this physics problem into Lean declarations in the covered file, with theorem and lemma proof bodies written as `by sorry`.
\end{theorem}
\begin{proof}
This is an autoformalization task, not a proof task. Produce statements that can later be proved by the physics prover without weakening the physical model.
\end{proof}

% --- Archon declaration topology begin ---
\section{Declaration topology}
\begin{definition}[acceleration Dimension]
  \label{def:physics:phyx-mini-0891:phyxminiproblems-problemphyxmini0891-accelerationdimension}
  \lean{PhyXMiniProblems.ProblemPhyXMini0891.accelerationDimension}
  Dimension of acceleration, L T⁻²
\end{definition}
\begin{proof}
  This is the declared model definition of acceleration Dimension.
\end{proof}
\begin{definition}[force Dimension]
  \label{def:physics:phyx-mini-0891:phyxminiproblems-problemphyxmini0891-forcedimension}
  \lean{PhyXMiniProblems.ProblemPhyXMini0891.forceDimension}
  Dimension of force, M L T⁻²
\end{definition}
\begin{proof}
  This is the declared model definition of force Dimension.
\end{proof}
\begin{definition}[spring Constant Dimension]
  \label{def:physics:phyx-mini-0891:phyxminiproblems-problemphyxmini0891-springconstantdimension}
  \lean{PhyXMiniProblems.ProblemPhyXMini0891.springConstantDimension}
  Dimension of a linear spring constant, M T⁻² = N/m
\end{definition}
\begin{proof}
  This is the declared model definition of spring Constant Dimension.
\end{proof}
\begin{definition}[electric Field Strength Dimension]
  \label{def:physics:phyx-mini-0891:phyxminiproblems-problemphyxmini0891-electricfieldstrengthdimension}
  \lean{PhyXMiniProblems.ProblemPhyXMini0891.electricFieldStrengthDimension}
  Dimension of electric-field strength, M L T⁻² C⁻¹ = N/C
\end{definition}
\begin{proof}
  This is the declared model definition of electric Field Strength Dimension.
\end{proof}
\begin{definition}[Mass Quantity]
  \label{def:physics:phyx-mini-0891:phyxminiproblems-problemphyxmini0891-massquantity}
  \lean{PhyXMiniProblems.ProblemPhyXMini0891.MassQuantity}
  A nonnegative, unit-independent physical mass
\end{definition}
\begin{proof}
  This is the declared model definition of Mass Quantity.
\end{proof}
\begin{definition}[Length Quantity]
  \label{def:physics:phyx-mini-0891:phyxminiproblems-problemphyxmini0891-lengthquantity}
  \lean{PhyXMiniProblems.ProblemPhyXMini0891.LengthQuantity}
  A nonnegative, unit-independent physical length
\end{definition}
\begin{proof}
  This is the declared model definition of Length Quantity.
\end{proof}
\begin{definition}[Charge Magnitude Quantity]
  \label{def:physics:phyx-mini-0891:phyxminiproblems-problemphyxmini0891-chargemagnitudequantity}
  \lean{PhyXMiniProblems.ProblemPhyXMini0891.ChargeMagnitudeQuantity}
  A nonnegative electric-charge magnitude
\end{definition}
\begin{proof}
  This is the declared model definition of Charge Magnitude Quantity.
\end{proof}
\begin{definition}[Acceleration Magnitude Quantity]
  \label{def:physics:phyx-mini-0891:phyxminiproblems-problemphyxmini0891-accelerationmagnitudequantity}
  \lean{PhyXMiniProblems.ProblemPhyXMini0891.AccelerationMagnitudeQuantity}
  \uses{def:physics:phyx-mini-0891:phyxminiproblems-problemphyxmini0891-accelerationdimension}
  A nonnegative acceleration magnitude
\end{definition}
\begin{proof}
  This is the declared model definition of Acceleration Magnitude Quantity.
\end{proof}
\begin{definition}[Force Magnitude Quantity]
  \label{def:physics:phyx-mini-0891:phyxminiproblems-problemphyxmini0891-forcemagnitudequantity}
  \lean{PhyXMiniProblems.ProblemPhyXMini0891.ForceMagnitudeQuantity}
  \uses{def:physics:phyx-mini-0891:phyxminiproblems-problemphyxmini0891-forcedimension}
  A nonnegative force magnitude
\end{definition}
\begin{proof}
  This is the declared model definition of Force Magnitude Quantity.
\end{proof}
\begin{definition}[Spring Constant Quantity]
  \label{def:physics:phyx-mini-0891:phyxminiproblems-problemphyxmini0891-springconstantquantity}
  \lean{PhyXMiniProblems.ProblemPhyXMini0891.SpringConstantQuantity}
  \uses{def:physics:phyx-mini-0891:phyxminiproblems-problemphyxmini0891-springconstantdimension}
  A nonnegative linear spring constant
\end{definition}
\begin{proof}
  This is the declared model definition of Spring Constant Quantity.
\end{proof}
\begin{definition}[Electric Field Strength Quantity]
  \label{def:physics:phyx-mini-0891:phyxminiproblems-problemphyxmini0891-electricfieldstrengthquantity}
  \lean{PhyXMiniProblems.ProblemPhyXMini0891.ElectricFieldStrengthQuantity}
  \uses{def:physics:phyx-mini-0891:phyxminiproblems-problemphyxmini0891-electricfieldstrengthdimension}
  A nonnegative electric-field strength
\end{definition}
\begin{proof}
  This is the declared model definition of Electric Field Strength Quantity.
\end{proof}
\begin{definition}[mass In Kilograms]
  \label{def:physics:phyx-mini-0891:phyxminiproblems-problemphyxmini0891-massinkilograms}
  \lean{PhyXMiniProblems.ProblemPhyXMini0891.massInKilograms}
  \uses{def:physics:phyx-mini-0891:phyxminiproblems-problemphyxmini0891-massquantity}
  Coherent-SI mass readout in kilograms
\end{definition}
\begin{proof}
  This is the declared model definition of mass In Kilograms.
\end{proof}
\begin{definition}[mass In Grams]
  \label{def:physics:phyx-mini-0891:phyxminiproblems-problemphyxmini0891-massingrams}
  \lean{PhyXMiniProblems.ProblemPhyXMini0891.massInGrams}
  \uses{def:physics:phyx-mini-0891:phyxminiproblems-problemphyxmini0891-massquantity, def:physics:phyx-mini-0891:phyxminiproblems-problemphyxmini0891-massinkilograms}
  Figure-scale mass readout in grams
\end{definition}
\begin{proof}
  This is the declared model definition of mass In Grams.
\end{proof}
\begin{definition}[length In Meters]
  \label{def:physics:phyx-mini-0891:phyxminiproblems-problemphyxmini0891-lengthinmeters}
  \lean{PhyXMiniProblems.ProblemPhyXMini0891.lengthInMeters}
  \uses{def:physics:phyx-mini-0891:phyxminiproblems-problemphyxmini0891-lengthquantity}
  Coherent-SI length readout in metres
\end{definition}
\begin{proof}
  This is the declared model definition of length In Meters.
\end{proof}
\begin{definition}[length In Centimeters]
  \label{def:physics:phyx-mini-0891:phyxminiproblems-problemphyxmini0891-lengthincentimeters}
  \lean{PhyXMiniProblems.ProblemPhyXMini0891.lengthInCentimeters}
  \uses{def:physics:phyx-mini-0891:phyxminiproblems-problemphyxmini0891-lengthquantity, def:physics:phyx-mini-0891:phyxminiproblems-problemphyxmini0891-lengthinmeters}
  Figure-scale length readout in centimetres
\end{definition}
\begin{proof}
  This is the declared model definition of length In Centimeters.
\end{proof}
\begin{definition}[charge In Coulombs]
  \label{def:physics:phyx-mini-0891:phyxminiproblems-problemphyxmini0891-chargeincoulombs}
  \lean{PhyXMiniProblems.ProblemPhyXMini0891.chargeInCoulombs}
  \uses{def:physics:phyx-mini-0891:phyxminiproblems-problemphyxmini0891-chargemagnitudequantity}
  Coherent-SI charge-magnitude readout in coulombs
\end{definition}
\begin{proof}
  This is the declared model definition of charge In Coulombs.
\end{proof}
\begin{definition}[charge In Nanocoulombs]
  \label{def:physics:phyx-mini-0891:phyxminiproblems-problemphyxmini0891-chargeinnanocoulombs}
  \lean{PhyXMiniProblems.ProblemPhyXMini0891.chargeInNanocoulombs}
  \uses{def:physics:phyx-mini-0891:phyxminiproblems-problemphyxmini0891-chargemagnitudequantity, def:physics:phyx-mini-0891:phyxminiproblems-problemphyxmini0891-chargeincoulombs}
  Figure-scale charge-magnitude readout in nanocoulombs
\end{definition}
\begin{proof}
  This is the declared model definition of charge In Nanocoulombs.
\end{proof}
\begin{definition}[acceleration In Meters Per Second Squared]
  \label{def:physics:phyx-mini-0891:phyxminiproblems-problemphyxmini0891-accelerationinmeterspersecondsquared}
  \lean{PhyXMiniProblems.ProblemPhyXMini0891.accelerationInMetersPerSecondSquared}
  \uses{def:physics:phyx-mini-0891:phyxminiproblems-problemphyxmini0891-accelerationmagnitudequantity}
  Coherent-SI acceleration readout in metres per second squared
\end{definition}
\begin{proof}
  This is the declared model definition of acceleration In Meters Per Second Squared.
\end{proof}
\begin{definition}[force In Newtons]
  \label{def:physics:phyx-mini-0891:phyxminiproblems-problemphyxmini0891-forceinnewtons}
  \lean{PhyXMiniProblems.ProblemPhyXMini0891.forceInNewtons}
  \uses{def:physics:phyx-mini-0891:phyxminiproblems-problemphyxmini0891-forcemagnitudequantity}
  Coherent-SI force-magnitude readout in newtons
\end{definition}
\begin{proof}
  This is the declared model definition of force In Newtons.
\end{proof}
\begin{definition}[spring Constant In Newtons Per Meter]
  \label{def:physics:phyx-mini-0891:phyxminiproblems-problemphyxmini0891-springconstantinnewtonspermeter}
  \lean{PhyXMiniProblems.ProblemPhyXMini0891.springConstantInNewtonsPerMeter}
  \uses{def:physics:phyx-mini-0891:phyxminiproblems-problemphyxmini0891-springconstantquantity}
  Coherent-SI spring-constant readout in newtons per metre
\end{definition}
\begin{proof}
  This is the declared model definition of spring Constant In Newtons Per Meter.
\end{proof}
\begin{definition}[electric Field Strength In Newtons Per Coulomb]
  \label{def:physics:phyx-mini-0891:phyxminiproblems-problemphyxmini0891-electricfieldstrengthinnewtonspercoulomb}
  \lean{PhyXMiniProblems.ProblemPhyXMini0891.electricFieldStrengthInNewtonsPerCoulomb}
  \uses{def:physics:phyx-mini-0891:phyxminiproblems-problemphyxmini0891-electricfieldstrengthquantity}
  Coherent-SI electric-field-strength readout in newtons per coulomb
\end{definition}
\begin{proof}
  This is the declared model definition of electric Field Strength In Newtons Per Coulomb.
\end{proof}
\begin{definition}[degrees To Radians]
  \label{def:physics:phyx-mini-0891:phyxminiproblems-problemphyxmini0891-degreestoradians}
  \lean{PhyXMiniProblems.ProblemPhyXMini0891.degreesToRadians}
  Convert a literal degree readout into radians
\end{definition}
\begin{proof}
  This is the declared model definition of degrees To Radians.
\end{proof}
\begin{definition}[Spring Material]
  \label{def:physics:phyx-mini-0891:phyxminiproblems-problemphyxmini0891-springmaterial}
  \lean{PhyXMiniProblems.ProblemPhyXMini0891.SpringMaterial}
  Constitutive/electrical classification of the spring
\end{definition}
\begin{proof}
  This is the declared model definition of Spring Material.
\end{proof}
\begin{definition}[Ball Force Role]
  \label{def:physics:phyx-mini-0891:phyxminiproblems-problemphyxmini0891-ballforcerole}
  \lean{PhyXMiniProblems.ProblemPhyXMini0891.BallForceRole}
  Force roles in the free-body diagram of the suspended ball
\end{definition}
\begin{proof}
  This is the declared model definition of Ball Force Role.
\end{proof}
\begin{definition}[Planar Direction]
  \label{def:physics:phyx-mini-0891:phyxminiproblems-problemphyxmini0891-planardirection}
  \lean{PhyXMiniProblems.ProblemPhyXMini0891.PlanarDirection}
  Qualitative planar directions needed by the free-body diagram
\end{definition}
\begin{proof}
  This is the declared model definition of Planar Direction.
\end{proof}
\begin{definition}[Charged Ball Spring Figure]
  \label{def:physics:phyx-mini-0891:phyxminiproblems-problemphyxmini0891-chargedballspringfigure}
  \lean{PhyXMiniProblems.ProblemPhyXMini0891.ChargedBallSpringFigure}
  Literal readouts and visible features of image 891.png. The 60 cm entry is named as a tether label because that is where it is placed in the raster
\end{definition}
\begin{proof}
  This is the declared model definition of Charged Ball Spring Figure.
\end{proof}
\begin{definition}[Charged Ball Spring Setup]
  \label{def:physics:phyx-mini-0891:phyxminiproblems-problemphyxmini0891-chargedballspringsetup}
  \lean{PhyXMiniProblems.ProblemPhyXMini0891.ChargedBallSpringSetup}
  \uses{def:physics:phyx-mini-0891:phyxminiproblems-problemphyxmini0891-massquantity, def:physics:phyx-mini-0891:phyxminiproblems-problemphyxmini0891-lengthquantity, def:physics:phyx-mini-0891:phyxminiproblems-problemphyxmini0891-chargemagnitudequantity, def:physics:phyx-mini-0891:phyxminiproblems-problemphyxmini0891-accelerationmagnitudequantity, def:physics:phyx-mini-0891:phyxminiproblems-problemphyxmini0891-forcemagnitudequantity, def:physics:phyx-mini-0891:phyxminiproblems-problemphyxmini0891-springconstantquantity, def:physics:phyx-mini-0891:phyxminiproblems-problemphyxmini0891-electricfieldstrengthquantity, def:physics:phyx-mini-0891:phyxminiproblems-problemphyxmini0891-springmaterial, def:physics:phyx-mini-0891:phyxminiproblems-problemphyxmini0891-ballforcerole, def:physics:phyx-mini-0891:phyxminiproblems-problemphyxmini0891-planardirection, def:physics:phyx-mini-0891:phyxminiproblems-problemphyxmini0891-chargedballspringfigure}
  Independent physical quantities of the apparatus. In particular, electricFieldStrength is not defined from an answer choice. Physlib's full spacetime electric field is retained in addition to its nonnegative SI strength readout
\end{definition}
\begin{proof}
  This is the declared model definition of Charged Ball Spring Setup.
\end{proof}
\begin{definition}[Matches Charged Ball Spring Figure]
  \label{def:physics:phyx-mini-0891:phyxminiproblems-problemphyxmini0891-matcheschargedballspringfigure}
  \lean{PhyXMiniProblems.ProblemPhyXMini0891.MatchesChargedBallSpringFigure}
  \uses{def:physics:phyx-mini-0891:phyxminiproblems-problemphyxmini0891-massingrams, def:physics:phyx-mini-0891:phyxminiproblems-problemphyxmini0891-lengthincentimeters, def:physics:phyx-mini-0891:phyxminiproblems-problemphyxmini0891-chargeinnanocoulombs, def:physics:phyx-mini-0891:phyxminiproblems-problemphyxmini0891-springconstantinnewtonspermeter, def:physics:phyx-mini-0891:phyxminiproblems-problemphyxmini0891-degreestoradians, def:physics:phyx-mini-0891:phyxminiproblems-problemphyxmini0891-chargedballspringsetup}
  Visible features and numerical labels transcribed from the primary image
\end{definition}
\begin{proof}
  This is the declared model definition of Matches Charged Ball Spring Figure.
\end{proof}
\begin{definition}[Matches Charged Ball Spring Scenario]
  \label{def:physics:phyx-mini-0891:phyxminiproblems-problemphyxmini0891-matcheschargedballspringscenario}
  \lean{PhyXMiniProblems.ProblemPhyXMini0891.MatchesChargedBallSpringScenario}
  \uses{def:physics:phyx-mini-0891:phyxminiproblems-problemphyxmini0891-chargedballspringsetup}
  The qualitative scenario, including the insulating role of the plastic spring and the directions of all four forces on the ball. Charge preservation is a physical consequence of the stated insulation, not an electric-field answer
\end{definition}
\begin{proof}
  This is the declared model definition of Matches Charged Ball Spring Scenario.
\end{proof}
\begin{definition}[Uses Standard Terrestrial Gravity]
  \label{def:physics:phyx-mini-0891:phyxminiproblems-problemphyxmini0891-usesstandardterrestrialgravity}
  \lean{PhyXMiniProblems.ProblemPhyXMini0891.UsesStandardTerrestrialGravity}
  \uses{def:physics:phyx-mini-0891:phyxminiproblems-problemphyxmini0891-accelerationinmeterspersecondsquared, def:physics:phyx-mini-0891:phyxminiproblems-problemphyxmini0891-chargedballspringsetup}
  Standard terrestrial gravitational acceleration used by the calculation
\end{definition}
\begin{proof}
  This is the declared model definition of Uses Standard Terrestrial Gravity.
\end{proof}
\begin{definition}[Has Physical Charged Ball Spring Parameters]
  \label{def:physics:phyx-mini-0891:phyxminiproblems-problemphyxmini0891-hasphysicalchargedballspringparameters}
  \lean{PhyXMiniProblems.ProblemPhyXMini0891.HasPhysicalChargedBallSpringParameters}
  \uses{def:physics:phyx-mini-0891:phyxminiproblems-problemphyxmini0891-massinkilograms, def:physics:phyx-mini-0891:phyxminiproblems-problemphyxmini0891-lengthinmeters, def:physics:phyx-mini-0891:phyxminiproblems-problemphyxmini0891-chargeincoulombs, def:physics:phyx-mini-0891:phyxminiproblems-problemphyxmini0891-accelerationinmeterspersecondsquared, def:physics:phyx-mini-0891:phyxminiproblems-problemphyxmini0891-forceinnewtons, def:physics:phyx-mini-0891:phyxminiproblems-problemphyxmini0891-springconstantinnewtonspermeter, def:physics:phyx-mini-0891:phyxminiproblems-problemphyxmini0891-electricfieldstrengthinnewtonspercoulomb, def:physics:phyx-mini-0891:phyxminiproblems-problemphyxmini0891-chargedballspringsetup}
  Positivity and the acute-angle regime required by the physical setup
\end{definition}
\begin{proof}
  This is the declared model definition of Has Physical Charged Ball Spring Parameters.
\end{proof}
\begin{definition}[Satisfies Deflected Tether And Spring Geometry]
  \label{def:physics:phyx-mini-0891:phyxminiproblems-problemphyxmini0891-satisfiesdeflectedtetherandspringgeometry}
  \lean{PhyXMiniProblems.ProblemPhyXMini0891.SatisfiesDeflectedTetherAndSpringGeometry}
  \uses{def:physics:phyx-mini-0891:phyxminiproblems-problemphyxmini0891-lengthinmeters, def:physics:phyx-mini-0891:phyxminiproblems-problemphyxmini0891-chargedballspringsetup}
  The dashed line is the zero-extension reference. Consequently the ball's horizontal displacement is L sin θ, and that displacement is exactly the extension of the horizontal spring
\end{definition}
\begin{proof}
  This is the declared model definition of Satisfies Deflected Tether And Spring Geometry.
\end{proof}
\begin{definition}[Satisfies Uniform Horizontal Electric Field]
  \label{def:physics:phyx-mini-0891:phyxminiproblems-problemphyxmini0891-satisfiesuniformhorizontalelectricfield}
  \lean{PhyXMiniProblems.ProblemPhyXMini0891.SatisfiesUniformHorizontalElectricField}
  \uses{def:physics:phyx-mini-0891:phyxminiproblems-problemphyxmini0891-electricfieldstrengthinnewtonspercoulomb, def:physics:phyx-mini-0891:phyxminiproblems-problemphyxmini0891-chargedballspringsetup}
  The spacetime field is uniform and equal to the displayed rightward vector. Coordinates 0, 1, and 2 are respectively rightward horizontal, upward vertical, and out of the page
\end{definition}
\begin{proof}
  This is the declared model definition of Satisfies Uniform Horizontal Electric Field.
\end{proof}
\begin{definition}[Satisfies Charged Ball Static Equilibrium]
  \label{def:physics:phyx-mini-0891:phyxminiproblems-problemphyxmini0891-satisfieschargedballstaticequilibrium}
  \lean{PhyXMiniProblems.ProblemPhyXMini0891.SatisfiesChargedBallStaticEquilibrium}
  \uses{def:physics:phyx-mini-0891:phyxminiproblems-problemphyxmini0891-massinkilograms, def:physics:phyx-mini-0891:phyxminiproblems-problemphyxmini0891-lengthinmeters, def:physics:phyx-mini-0891:phyxminiproblems-problemphyxmini0891-chargeincoulombs, def:physics:phyx-mini-0891:phyxminiproblems-problemphyxmini0891-accelerationinmeterspersecondsquared, def:physics:phyx-mini-0891:phyxminiproblems-problemphyxmini0891-forceinnewtons, def:physics:phyx-mini-0891:phyxminiproblems-problemphyxmini0891-springconstantinnewtonspermeter, def:physics:phyx-mini-0891:phyxminiproblems-problemphyxmini0891-electricfieldstrengthinnewtonspercoulomb, def:physics:phyx-mini-0891:phyxminiproblems-problemphyxmini0891-chargedballspringsetup}
  Hooke's law, electric force, weight, and static force balance, all stated at a coherent-SI boundary. The horizontal balance has the rightward electric force opposed by both the spring and the tether's horizontal component. No numerical field strength or answer choice occurs in these laws
\end{definition}
\begin{proof}
  This is the declared model definition of Satisfies Charged Ball Static Equilibrium.
\end{proof}
\begin{theorem}[electric Field Strength closed Form]
  \label{thm:physics:phyx-mini-0891:phyxminiproblems-problemphyxmini0891-electricfieldstrength-closedform}
  \lean{PhyXMiniProblems.ProblemPhyXMini0891.electricFieldStrength_closedForm}
  \uses{def:physics:phyx-mini-0891:phyxminiproblems-problemphyxmini0891-massinkilograms, def:physics:phyx-mini-0891:phyxminiproblems-problemphyxmini0891-lengthinmeters, def:physics:phyx-mini-0891:phyxminiproblems-problemphyxmini0891-chargeincoulombs, def:physics:phyx-mini-0891:phyxminiproblems-problemphyxmini0891-accelerationinmeterspersecondsquared, def:physics:phyx-mini-0891:phyxminiproblems-problemphyxmini0891-springconstantinnewtonspermeter, def:physics:phyx-mini-0891:phyxminiproblems-problemphyxmini0891-electricfieldstrengthinnewtonspercoulomb, def:physics:phyx-mini-0891:phyxminiproblems-problemphyxmini0891-chargedballspringsetup, def:physics:phyx-mini-0891:phyxminiproblems-problemphyxmini0891-matcheschargedballspringfigure, def:physics:phyx-mini-0891:phyxminiproblems-problemphyxmini0891-matcheschargedballspringscenario, def:physics:phyx-mini-0891:phyxminiproblems-problemphyxmini0891-usesstandardterrestrialgravity, def:physics:phyx-mini-0891:phyxminiproblems-problemphyxmini0891-hasphysicalchargedballspringparameters, def:physics:phyx-mini-0891:phyxminiproblems-problemphyxmini0891-satisfiesdeflectedtetherandspringgeometry, def:physics:phyx-mini-0891:phyxminiproblems-problemphyxmini0891-satisfiesuniformhorizontalelectricfield, def:physics:phyx-mini-0891:phyxminiproblems-problemphyxmini0891-satisfieschargedballstaticequilibrium}
  Eliminating tension and the two intermediate force magnitudes gives the exact static-equilibrium field-strength formula. It is a conclusion, not a premise or a definition of electricFieldStrength
\end{theorem}
\begin{proof}
  The conclusion follows from the governing relations and figure readouts named in the statement of electric Field Strength closed Form.
\end{proof}
\begin{definition}[Answer Choice]
  \label{def:physics:phyx-mini-0891:phyxminiproblems-problemphyxmini0891-answerchoice}
  \lean{PhyXMiniProblems.ProblemPhyXMini0891.AnswerChoice}
  Labels of the four electric-field-strength choices in the source
\end{definition}
\begin{proof}
  This is the declared model definition of Answer Choice.
\end{proof}
\begin{definition}[answer Choice Field Strength In Newtons Per Coulomb]
  \label{def:physics:phyx-mini-0891:phyxminiproblems-problemphyxmini0891-answerchoicefieldstrengthinnewtonspercoulomb}
  \lean{PhyXMiniProblems.ProblemPhyXMini0891.answerChoiceFieldStrengthInNewtonsPerCoulomb}
  \uses{def:physics:phyx-mini-0891:phyxminiproblems-problemphyxmini0891-answerchoice}
  Literal field-strength readout, in newtons per coulomb, for each choice
\end{definition}
\begin{proof}
  This is the declared model definition of answer Choice Field Strength In Newtons Per Coulomb.
\end{proof}
\begin{definition}[recorded Dataset Answer]
  \label{def:physics:phyx-mini-0891:phyxminiproblems-problemphyxmini0891-recordeddatasetanswer}
  \lean{PhyXMiniProblems.ProblemPhyXMini0891.recordedDatasetAnswer}
  \uses{def:physics:phyx-mini-0891:phyxminiproblems-problemphyxmini0891-answerchoice}
  The answer label recorded by the source dataset
\end{definition}
\begin{proof}
  This is the declared model definition of recorded Dataset Answer.
\end{proof}
\begin{definition}[answer Choice Absolute Error]
  \label{def:physics:phyx-mini-0891:phyxminiproblems-problemphyxmini0891-answerchoiceabsoluteerror}
  \lean{PhyXMiniProblems.ProblemPhyXMini0891.answerChoiceAbsoluteError}
  \uses{def:physics:phyx-mini-0891:phyxminiproblems-problemphyxmini0891-electricfieldstrengthinnewtonspercoulomb, def:physics:phyx-mini-0891:phyxminiproblems-problemphyxmini0891-chargedballspringsetup, def:physics:phyx-mini-0891:phyxminiproblems-problemphyxmini0891-answerchoice, def:physics:phyx-mini-0891:phyxminiproblems-problemphyxmini0891-answerchoicefieldstrengthinnewtonspercoulomb}
  Absolute error of a displayed choice from the modeled physical field
\end{definition}
\begin{proof}
  This is the declared model definition of answer Choice Absolute Error.
\end{proof}
\begin{definition}[Is Unique Closest Answer]
  \label{def:physics:phyx-mini-0891:phyxminiproblems-problemphyxmini0891-isuniqueclosestanswer}
  \lean{PhyXMiniProblems.ProblemPhyXMini0891.IsUniqueClosestAnswer}
  \uses{def:physics:phyx-mini-0891:phyxminiproblems-problemphyxmini0891-chargedballspringsetup, def:physics:phyx-mini-0891:phyxminiproblems-problemphyxmini0891-answerchoice, def:physics:phyx-mini-0891:phyxminiproblems-problemphyxmini0891-answerchoiceabsoluteerror}
  A selected choice is strictly closer than every other displayed choice
\end{definition}
\begin{proof}
  This is the declared model definition of Is Unique Closest Answer.
\end{proof}
% --- Archon declaration topology end ---
% --- Archon physics formalization source end ---
